% Hockey Pool Analytics Pipeline — Full Documentation
% Compile with: pdflatex hockey_pool_pipeline.tex
% All packages are standard TeX Live / MiKTeX; no external dependencies.

\documentclass[11pt, a4paper]{article}

\usepackage[margin=2.5cm]{geometry}
\usepackage{amsmath, amssymb}
\usepackage{booktabs}
\usepackage{array}
\usepackage{longtable}
\usepackage[hidelinks]{hyperref}
\usepackage{listings}
\usepackage{xcolor}
\usepackage{graphicx}
\usepackage{parskip}
\usepackage{enumitem}
\usepackage{microtype}
\usepackage{caption}

% ---- code listing style ----
\definecolor{codebg}{rgb}{0.97,0.97,0.97}
\definecolor{comment}{rgb}{0.4,0.4,0.4}
\lstset{
  backgroundcolor=\color{codebg},
  basicstyle=\ttfamily\small,
  commentstyle=\color{comment},
  keywordstyle=\bfseries,
  frame=single,
  framerule=0.4pt,
  breaklines=true,
  captionpos=b,
  aboveskip=8pt,
  belowskip=4pt,
}

% ---- table column types ----
\newcolumntype{L}[1]{>{\raggedright\arraybackslash}p{#1}}
\newcolumntype{C}[1]{>{\centering\arraybackslash}p{#1}}

\title{\textbf{Hockey Pool Analytics Pipeline}\\[0.4em]
       \large Architecture, Models, Diagnostics, and Results}
\author{Julian di Giovanni}
\date{2026-07-24}

\begin{document}

\maketitle
\tableofcontents
\newpage

% ===========================================================================
\section{Introduction}
\label{sec:intro}
% ===========================================================================

\subsection{Project purpose}

This project produces draft-day player rankings for a custom fantasy hockey pool held once per
year. The pipeline is rebuilt annually from a shared codebase with minimal configuration changes,
following the principle: \emph{copy last season's config, update a few fields, run \texttt{--stage all}}.

The output is an Excel workbook and set of CSVs ranking every NHL player by predicted pool points
for the upcoming season, broken down by position (forwards, defense, goalies) and combined into
an overall ranking.

\subsection{Pool scoring rules}

Scoring rules are encoded in \texttt{<season>/config.yaml} under the \texttt{scoring:} key and
are the same year-to-year. Table~\ref{tab:scoring} summarises the current rules.

\begin{table}[h]
\centering
\caption{Pool scoring rules (2026-27 season)}
\label{tab:scoring}
\begin{tabular}{llr}
\toprule
Position & Event & Points \\
\midrule
Forwards  & Goal                             & 1 \\
          & Assist                           & 1 \\
\midrule
Defense   & Goal                             & 1 \\
          & Assist                           & 1 \\
          & Plus/minus                       & 1 per unit \\
\midrule
Goalies   & Win (non-shutout)                & 1 \\
          & Shutout win (total, not additive)& 3 \\
          & Best GAA among $\geq$40 GP qualifiers & +10 \\
          & Best save\% among $\geq$40 GP qualifiers & +10 \\
\bottomrule
\end{tabular}
\end{table}

\subsection{Annual workflow}

\begin{enumerate}[leftmargin=*]
  \item \textbf{July} (off-season): run \texttt{--stage all} to refresh historical data through
        the just-completed season and generate 2026-27 predictions.
  \item \textbf{Draft day}: use the output rankings and XLSX workbook to make pick decisions.
  \item \textbf{Following July}: copy \texttt{config.yaml} to the new season folder, update
        \texttt{season}, \texttt{season\_id}, \texttt{roster\_as\_of\_date}, and any
        \texttt{trade\_overrides}. Run \texttt{--stage all}.
\end{enumerate}


% ===========================================================================
\section{System Architecture}
\label{sec:architecture}
% ===========================================================================

\subsection{Folder layout}

\begin{lstlisting}[language=bash, caption={Repository structure}]
common/               shared library, reused every season
  config.py             SeasonConfig dataclass (loads config.yaml)
  scrape/
    nhl_api.py            NHL API client (stats, rosters, game logs)
    rosters.py            current-team lookup (handles off-season moves)
    sources/
      moneypuck.py          MoneyPuck CSV scraper (xG, shot quality, GSAx)
  clean/
    merge.py              merges skater/goalie/team data into model CSVs
  features/
    engineering.py        lag features, EB shrinkage, GSAx, leakage guard
  models/
    training.py           5-algorithm zoo, CV, feature selection, OOS eval
    forwards.py           forward target + feature config
    defense.py            defense target + feature config
    goalies.py            goalie target + feature config
    pool_ranking.py       scoring, blending, ranking, output writing
  diagnostics/
    reconcile.py          cross-source reconciliation (NHL API vs MoneyPuck)
    model_report.py       metrics table, PvA plots, rolling OOS plot
  pipeline.py             stage orchestration (scrape/clean/train/predict)

run_season.py         CLI entry point

<season>/             one directory per draft year, e.g. 2026-27/
  config.yaml           season parameters and scoring rules
  data/{raw,processed}/ scraped + cleaned data (gitignored)
  models/               persisted joblib artifacts (gitignored)
  results/              rankings, XLSX, report, diagnostics CSV (gitignored)
  plots/                PvA and rolling-OOS plots (gitignored)
\end{lstlisting}

\subsection{Pipeline stages}

The four stages run in sequence via:

\begin{lstlisting}[language=bash]
python3 run_season.py --season 2026-27 --stage all
# or individually:
python3 run_season.py --season 2026-27 --stage scrape
python3 run_season.py --season 2026-27 --stage clean
python3 run_season.py --season 2026-27 --stage train
python3 run_season.py --season 2026-27 --stage predict
\end{lstlisting}

\begin{description}[leftmargin=2cm, style=nextline]
  \item[\texttt{scrape}] Pull skater/goalie stats, team stats, rosters, and game logs from the
        NHL API; pull xG/shot-quality/GSAx data from MoneyPuck. Seed from prior season's raw
        directory to avoid re-downloading unchanged historical seasons.
  \item[\texttt{clean}] Merge position stats with team context; compute per-game rates; output
        \texttt{skater\_team\_data.csv} and \texttt{goalie\_team\_data.csv}.
  \item[\texttt{train}] Run feature engineering, train the 5-algorithm zoo for each target,
        select the best model, persist to \texttt{models/<target>.joblib}.
  \item[\texttt{predict}] Load persisted models, append prediction rows for the target season,
        score pool points, write rankings, generate diagnostics. Runs in $\sim$2s (no retraining).
\end{description}


% ===========================================================================
\section{Data Sources}
\label{sec:data}
% ===========================================================================

\subsection{NHL API (primary)}

\textbf{Endpoint:} \texttt{api.nhle.com} and \texttt{api-web.nhle.com}.
Free, no authentication key required. Historical data from the 2008-09 season onward.

Data pulled per season: skater counting stats (goals, assists, shots, TOI, hits, blocks, PIM,
plus/minus), goalie stats (wins, losses, OTL, shutouts, saves, shots against, GAA, save\%),
team-level stats (goals for/against, points, games played), and current-season per-game logs.

Roster snapshots are pulled as of \texttt{roster\_as\_of\_date} in the season config, so
off-season trades, free-agent signings, and waiver moves are reflected automatically.

\subsection{MoneyPuck (secondary)}

\textbf{Source:} \texttt{moneypuck.com/data.htm} — free CSV downloads, no key required.
All 18 seasons (2008--2025) successfully fetched: 81,355 skater rows, 8,510 goalie rows.

MoneyPuck provides metrics unavailable from the NHL API:
\begin{itemize}[leftmargin=*]
  \item \textbf{Expected goals (xG)}: shot-quality-adjusted goal probability per shot
  \item \textbf{Danger-zone shot counts}: high / medium / low danger breakdowns
  \item \textbf{gameScore}: composite single-number performance index
  \item \textbf{GSAx} (goalie-specific): goals saved above expected
\end{itemize}

\textbf{Join strategy:} The integer \texttt{player\_id} is identical in both sources.
Join key: \texttt{player\_id} (integer) $+$ season start year. No name-matching is used.
Match rate: $\sim$94.5\%. The remaining $\sim$5.5\% are fringe players genuinely absent from
MoneyPuck's records (very few games played). A name-fallback join is attempted for ID misses
only.

\textbf{Rate-limit handling:} 1\,s between year fetches, 3\,s between player entities, 15\,s
retry on HTTP 429.

\subsection{Cross-source reconciliation}

A dedicated diagnostic script (\texttt{common/diagnostics/reconcile.py}) validates agreement
between NHL API counting stats and MoneyPuck. Results for the full historical dataset:
Pearson $r \geq 0.9999$ for goals, assists, points, and games\_played. The few
games\_played outliers in earlier seasons trace to historical surname collisions in MoneyPuck
(e.g.\ two players named ``Jones'' in the same season) — the player\_id join prevents these
from affecting the model.

\subsection{Sources not used}

\begin{description}[leftmargin=2cm, style=nextline]
  \item[Natural Stat Trick] Skipped. Free but requires a manually-approved access key before
        automated pulls. Its stats mostly overlap MoneyPuck's.
  \item[Evolving-Hockey] Skipped. The valuable parts (RAPM, GAR, projections, CSV downloads)
        are paywalled; the free tier adds nothing beyond MoneyPuck.
  \item[Elite Prospects] Skipped. Bio/prospect data; not stats-focused; scraping ToS is unclear.
\end{description}


% ===========================================================================
\section{Feature Engineering}
\label{sec:features}
% ===========================================================================

All feature engineering lives in \texttt{common/features/engineering.py}.

\subsection{Lag features}

The core predictor of next-season performance is last-season (and two-seasons-ago) performance.
All individual statistics, MoneyPuck metrics, and team-context columns are lagged before use:

\begin{equation}
  x_{i,t}^{\text{lag}k} = x_{i,t-k}, \quad k \in \{1, 2\}
\end{equation}

Lags are computed via a vectorized \texttt{groupby("player\_id").shift(k)}, producing columns
\texttt{<feature>\_lag1} and \texttt{<feature>\_lag2} for every eligible feature.
Rows missing all core lag columns (players with no prior history) are dropped from training.

\subsection{Leakage prevention}

The feature selector uses an \textbf{allowlist}, not a denylist. A column is eligible as a
model predictor only if it satisfies at least one of:

\begin{enumerate}[leftmargin=*]
  \item Contains \texttt{"lag"} in its name (i.e.\ it is a genuine lag column), or
  \item Contains \texttt{"covid"} in its name (COVID season indicators), or
  \item Contains any of the \texttt{\_ENGINEERED\_MARKERS} substrings:
        \texttt{\_career\_avg}, \texttt{\_hist\_avg}, \texttt{\_hist\_std}, \texttt{\_trend},
        \texttt{\_per\_game\_lag}, \texttt{rel\_team}, \texttt{teammate},
        \texttt{normalized\_team}, \texttt{\_performance}, \texttt{\_x\_}
\end{enumerate}

Contemporaneous current-season columns never pass this filter (unless explicitly listed and
engineered). This allowlist design catches leakage by construction — a new column added to the
data schema cannot slip through unless the author intentionally adds it to the markers list.

\subsection{Career averages}

For high-variance rate statistics, the single most-recent season lag is a noisy estimate. An
expanding career average provides a more stable signal:

\begin{equation}
  \bar{x}_{i,t}^{\text{career}} = \frac{1}{t-1} \sum_{s < t} x_{i,s}
\end{equation}

Implemented via \texttt{groupby("player\_id").shift(1).expanding().mean()} — the shift(1)
ensures only seasons prior to $t$ are included (no leakage). Produces columns named
\texttt{<feature>\_career\_avg}.

Career averages are computed for:
\begin{itemize}[leftmargin=*]
  \item \textbf{Goalies}: save\_pct, GAA, shutouts/game, wins/game, eb\_save\_pct, GSAx
  \item \textbf{Defense}: plus\_minus, plus\_minus\_per\_game
\end{itemize}

\subsection{Empirical Bayes save percentage}
\label{sec:eb}

Raw season save percentage is a binomial proportion subject to substantial sample-size noise.
A backup goalie who faced 300 shots and saved 275 (Sv\%\ = 0.917) has a much wider posterior
uncertainty than a starter who faced 1,800 shots at the same rate. The Empirical Bayes (EB)
approach (Thomas 2006; community standard in hockey analytics) shrinks each estimate toward the
league mean in proportion to the goalie's shot volume.

\textbf{Model:} Assume saves $\sim \text{Binomial}(n, p)$ where $n$ = shots faced and $p$ is
the true save rate. Place a Beta prior $p \sim \text{Beta}(\alpha, \beta)$.
The posterior mean is:

\begin{equation}
  \hat{p}_{i}^{\text{EB}} = \frac{\alpha + \text{saves}_i}{\alpha + \beta + n_i}
  \label{eq:eb}
\end{equation}

\textbf{Prior estimation:} $\alpha$ and $\beta$ are estimated from the data via method of
moments on all goalie-seasons with $n_i \geq 100$ shots:

\begin{align}
  \hat{\mu} &= \text{mean}(\hat{p}_i), \quad
  \hat{\sigma}^2 = \text{var}(\hat{p}_i) \\
  \hat{\alpha} &= \hat{\mu}\!\left(\frac{\hat{\mu}(1-\hat{\mu})}{\hat{\sigma}^2} - 1\right),
  \quad
  \hat{\beta} = \hat{\alpha}\,\frac{1 - \hat{\mu}}{\hat{\mu}}
\end{align}

If the pooled dataset contains fewer than 50 qualifying seasons, the fallback prior
$(\alpha, \beta) = (85.0,\; 8.5)$ is used, placing the prior mean at $\approx 0.909$ with
moderate concentration. The resulting column \texttt{eb\_save\_pct} is then lagged and
career-averaged before use.

\subsection{GSAx and shot-quality features}
\label{sec:gsax}

Raw GAA conflates shot volume (a team-defense quality) with goalie efficiency. MoneyPuck
provides expected goals against based on shot location and type, enabling a cleaner
individual-skill metric:

\begin{equation}
  \text{GSAx}_i = \text{xGoals\_against}_i - \text{goals\_against}_i
\end{equation}

A positive GSAx means the goalie \emph{outperformed} the shot difficulty they faced.
GSAx is more persistent season-to-season than raw save percentage (published YoY autocorrelation:
$r \approx 0.15$--$0.30$ for GSAx vs $r \approx 0.05$--$0.15$ for raw Sv\%).

Additionally, \texttt{hd\_shot\_pct} $= $ high-danger shots / total shots against captures the
shot-mix the defense allowed — a partial control for team quality independent of the goalie.

Both features are lagged and added to career averages.

\subsection{COVID and experience features}

\textbf{COVID indicators:} The 2019-20 season was suspended mid-year; the 2020-21 season was
shortened to 56 games. A binary \texttt{covid\_year} flag and interaction terms
(\texttt{<metric>\_x\_covid}) are added to allow models to down-weight these atypical seasons.

\textbf{Experience:} A polynomial in years played is added for each position:
\[
  \text{years\_played},\quad \text{years\_played}^2,\quad \text{years\_played}^3
\]
with a position-specific ``prime'' window (forwards: years 3--10; defense: years 3--12;
goalies: years 3--12) encoded as an additional binary indicator.

\subsection{Interaction features}

\texttt{engineer\_features()} generates pairwise lag$\times$lag interaction terms between a
small set of key metrics (e.g.\ plus\_minus\_per\_game $\times$ points\_per\_game for defense),
and lag $\times$ COVID interactions for the most volatile stats. These are allowlisted via the
\texttt{\_x\_} marker.


% ===========================================================================
\section{Modeling Framework}
\label{sec:modeling}
% ===========================================================================

\subsection{Algorithm zoo}

Five candidate algorithms are trained for each target. Table~\ref{tab:zoo} summarises the
hyperparameter search for each.

\begin{table}[h]
\centering
\caption{5-algorithm zoo and search configuration}
\label{tab:zoo}
\begin{tabular}{L{2.8cm} C{1.8cm} L{8cm}}
\toprule
Algorithm & Search & Hyperparameter grid \\
\midrule
Random Forest &
  Randomized (8) &
  \texttt{n\_estimators} $\in \{150,250\}$,
  \texttt{max\_depth} $\in \{12,18\}$,
  \texttt{min\_samples\_split} $\in \{15,25\}$,
  \texttt{min\_samples\_leaf} $\in \{8,12\}$,
  \texttt{max\_features} $\in \{\sqrt{p}, 0.4\}$ \\
\addlinespace
Gradient Boosting &
  Randomized (10) &
  \texttt{n\_estimators} $\in \{150,250\}$,
  \texttt{max\_depth} $\in \{3,5\}$,
  \texttt{learning\_rate} $\in \{0.05,0.1\}$,
  \texttt{subsample} $\in \{0.8,0.9\}$,
  early stopping: \texttt{validation\_fraction=0.2},
  \texttt{n\_iter\_no\_change=10} \\
\addlinespace
Ridge & Grid & \texttt{alpha} $\in \{10, 50, 100, 500\}$ \\
\addlinespace
Lasso & Grid &
  \texttt{alpha} $\in \{0.1, 0.5, 1.0, 5.0\}$,
  \texttt{max\_iter=3000} \\
\addlinespace
ElasticNet & Grid &
  \texttt{alpha} $\in \{0.1, 1.0, 5.0\}$,
  \texttt{l1\_ratio} $\in \{0.3, 0.5, 0.7\}$,
  \texttt{max\_iter=3000} \\
\bottomrule
\end{tabular}
\end{table}

Linear models (Ridge, Lasso, ElasticNet) receive a \texttt{RobustScaler} prior to fitting.
Tree models are scale-invariant and receive no scaling.

\subsection{Cross-validation strategy}

When the training data has a temporal ordering (season year), a \texttt{TimeSeriesSplit} is
used, holding out the most-recent block of the training set:

\begin{align}
  n_{\text{folds}} &= \min\!\left(5,\; \max\!\left(3,\; \lfloor n / 30 \rfloor\right)\right) \\
  s_{\text{test}} &= \min\!\left(0.25,\; \max\!\left(0.15,\; 40 / n\right)\right)
\end{align}

This ensures a minimum of 40 rows in the held-out test set while never exceeding 25\% of the
training data. Without temporal ordering, a shuffled \texttt{KFold} is used.

\subsection{Feature selection}

When the feature matrix has more than 20 columns, \texttt{SelectKBest} with F-statistic
univariate regression is applied before fitting:

\begin{equation}
  k = \min\!\left(50,\; \max\!\left(15,\; \lfloor p / 3 \rfloor\right)\right)
\end{equation}

where $p$ is the total number of eligible features. This reduces noise features while retaining
at least 15 predictors.

\subsection{Model selection score}

After fitting each algorithm, a \textbf{selection score} balances test-set performance with
overfitting:

\begin{equation}
  \text{score} = R^2_{\text{test}} - 0.1 \times \frac{R^2_{\text{train}} - R^2_{\text{test}}}{R^2_{\text{train}}}
  \label{eq:selection_score}
\end{equation}

A model with a high train $R^2$ but poor test $R^2$ is penalised. The algorithm with the
highest selection score is chosen as the production model for each target.

\subsection{Model persistence}

The winning model for each target is saved to \texttt{models/<target>.joblib} via
\texttt{joblib.dump}. The \texttt{--stage predict} step loads these artifacts ($\sim$2s) rather
than retraining ($\sim$60s). The \texttt{TrainedModel} dataclass stores:
\texttt{target}, \texttt{model\_type}, \texttt{estimator}, \texttt{selected\_features},
\texttt{scaler}, \texttt{metrics}, \texttt{best\_params}, \texttt{y\_test}, \texttt{y\_pred\_test}
(the latter two enabling diagnostic plots without retraining).


% ===========================================================================
\section{Out-of-Sample Validation}
\label{sec:oos}
% ===========================================================================

\subsection{Primary OOS method: train-all vs.\ exclude-latest}

Every target is validated via \texttt{train\_all\_vs\_exclude\_latest()}:

\begin{enumerate}[leftmargin=*]
  \item Train \textbf{model A} on all available seasons.
  \item Train \textbf{model B} on all seasons \emph{except} the most recent.
  \item Evaluate model B on the held-out most-recent season — this is the out-of-sample (OOS) evaluation.
  \item Choose the model with the higher selection score (Eq.~\ref{eq:selection_score}) for
        production prediction. In practice, model A (more data) wins most of the time.
\end{enumerate}

\subsection{OOS metrics}

For each model, the following metrics are reported for both the in-sample holdout and the OOS
evaluation:

\begin{description}[leftmargin=3cm, style=nextline]
  \item[$R^2$] Coefficient of determination. Fraction of variance explained.
        $R^2 < 0$ means the model is worse than predicting the training-set mean.
  \item[MAE] Mean absolute error. Scale-dependent; reported in same units as the target.
  \item[Baseline $R^2$] $R^2$ of a constant predictor (training-set mean). Any model with
        $R^2 < R^2_{\text{baseline}}$ is flagged \texttt{beats\_baseline=False}.
  \item[Spearman $\rho$] Rank correlation between predicted and actual values. Range $[-1, 1]$;
        0 = no ranking information. Directly measures the pool draft objective.
  \item[Concordance] Pairwise ranking accuracy:
        $C = \Pr[\hat{y}_i > \hat{y}_j \mid y_i > y_j]$.
        Computed as $({\tau}+1)/2$ where $\tau$ is Kendall's $\tau$.
        Range $[0,1]$; $C = 0.5$ = chance.
  \item[Directional accuracy] Fraction of players where the model correctly predicts whether
        they are above or below the mean: $\text{dir\_acc} = \Pr[\text{sign}(\hat{y}-\bar{\hat{y}}) = \text{sign}(y - \bar{y})]$.
  \item[Bias] Systematic over- or under-prediction: $\bar{\hat{y}} - \bar{y}$.
  \item[\texttt{top1\_match\_max}] Binary: does the model's argmax (predicted best player)
        match the actual best player in the OOS year? Analogously \texttt{top1\_match\_min}
        for the argmin (e.g.\ the goalie who would win the best-GAA bonus).
\end{description}

\textbf{Rationale for rank metrics:} $R^2 \approx 0$ does not mean the model has zero utility
for the pool. The draft is a ranking problem — what matters is whether the model correctly
orders players relative to each other, not the absolute predicted value. A model with
$R^2 = 0.05$ and concordance $= 0.65$ is genuinely useful.

\subsection{Rolling OOS evaluation}

An optional gold-standard time-series OOS is implemented in \texttt{rolling\_oos\_eval()}.
For each holdout year $t$ (requiring at least 5 prior training years), the model is trained on
all years $< t$ and evaluated on year $t$. This produces an $R^2$ curve over holdout years,
visualised in \texttt{plots/diagnostics/rolling\_oos\_all.png}. It is computationally expensive
(one full train per holdout year) and is triggered with the \texttt{--rolling-eval} flag.


% ===========================================================================
\section{Position-Specific Models}
\label{sec:positions}
% ===========================================================================

\subsection{Forwards}
\label{sec:forwards}

\textbf{Target:} \texttt{points\_per\_game} = (goals + assists) / games\_played.

\textbf{Algorithms:} All five (no restriction).

\textbf{Feature set:} All skater individual stats lagged 1 and 2 years; MoneyPuck xG and shot
metrics lagged; team context (goals for/against, power-play efficiency) lagged; experience
polynomial; COVID interactions. Forwards are split from defensemen before modelling.

\textbf{Observations:} All skaters with at least one full year of lag history in the training
data ($\sim$9{,}000--12{,}000 rows across 2010--2025).

\textbf{OOS performance (2026-07-24):} $R^2 = 0.665$, concordance $\approx 0.83$ — the healthiest
target in the pipeline. Points are moderately persistent: elite forwards remain elite.

\subsection{Defense}
\label{sec:defense}

\textbf{Targets:}
\begin{itemize}[leftmargin=*]
  \item \texttt{points\_per\_game} (goal + assists per game) — same construct as forwards.
  \item \texttt{plus\_minus\_per\_game} — net goal differential while on ice per game.
\end{itemize}

\textbf{Algorithms:} All five for both targets.

\textbf{Leakage:} \texttt{allow\_contemporaneous\_team=False} for both targets.
Contemporaneous team stats (e.g.\ this season's team goals-for) are NaN at prediction time.
Lagged team proxies (\texttt{team\_goals\_for\_per\_game\_lag1}, teammate-quality composites) are
used instead.

\textbf{Additional features:} Career averages for plus\_minus and plus\_minus\_per\_game, and
four lagged teammate-quality proxies (offensive strength, defensive strength, consistency,
elite-team indicator).

\textbf{Prediction blending:} Both defense targets are blended post-prediction with 2-year
historical per-game averages to correct for model compression at the extremes of the
distribution (Section~\ref{sec:blending}).

\textbf{OOS performance:} points/game $R^2 \approx 0.4$, concordance $\approx 0.72$;
plus\_minus/game $R^2 = 0.031$, concordance $\approx 0.54$.

\subsection{Goalies}
\label{sec:goalies}

\textbf{Targets:}

\begin{center}
\begin{tabular}{lll}
\toprule
Target & Internal name & Restriction \\
\midrule
Wins per game      & \texttt{wins\_per\_game}           & All goalies \\
Shutouts per game  & \texttt{shutouts\_per\_game}        & All goalies \\
Goals against avg  & \texttt{goals\_against\_average}    & $\geq$40 GP qualifiers only \\
Save percentage    & \texttt{save\_percentage}           & $\geq$40 GP qualifiers only \\
\bottomrule
\end{tabular}
\end{center}

\textbf{Algorithm restriction for GAA and save\%:} Only Ridge, Lasso, and ElasticNet
(\texttt{\_LINEAR\_ONLY}). Season-to-season autocorrelation for both stats is $r \approx 0.09$
(near zero). Tree models fit this noise and produce predictions that cluster around the training
mean; regularized linear models correctly regress toward the mean by design.

\textbf{Training filter:} All goalies (not just qualified ones) are included in training.
Previously, only goalies with $\geq 40$ GP were used, collapsing the training set from
$\sim$1{,}100 to $\sim$340 rows. The QUALIFIED\_ONLY distinction now applies only at prediction
time (bonus eligibility), not at training time.

\textbf{Key engineered features for goalies:}
\begin{itemize}[leftmargin=*]
  \item \texttt{eb\_save\_pct}: Empirical Bayes posterior mean (Section~\ref{sec:eb}) — more
        stable than raw save\_pct, especially for low-volume goalies.
  \item \texttt{gsax}: Goals saved above expected (Section~\ref{sec:gsax}) — lagged GSAx
        more persistent than raw Sv\%.
  \item \texttt{hd\_shot\_pct}: High-danger shot fraction — partial control for team defense quality.
  \item Career averages for all rate stats including \texttt{eb\_save\_pct} and \texttt{gsax}.
\end{itemize}

\textbf{OOS performance:} wins/game $R^2 \approx 0.10$; shutouts/game $R^2 \approx 0.05$;
GAA $R^2 = -0.203$; save\% $R^2 = -0.839$. The negative $R^2$ for precision stats is expected
given near-zero autocorrelation — this is a fundamental data property, not a model failure
(Thomas 2006). The bonus only needs correct argmax/argmin among qualified goalies, so
concordance $> 0.5$ is the relevant threshold.


% ===========================================================================
\section{Prediction Blending for Defense}
\label{sec:blending}
% ===========================================================================

\subsection{Problem: model compression}

Gradient boosting (and to a lesser extent other tree models) tends to compress predictions
toward the center of the training distribution. As a concrete example, Cale Makar's
points\_per\_game\_lag1 = 1.053 was predicted as 0.762 — the model ``averaged him toward the
pack.'' Separately, ElasticNet predicted near-constant plus\_minus for almost all defensemen
(heavy $\ell_1$ regularization zeroed most coefficients, leaving predictions clustered at
$\sim$0).

The result was that elite D-men ranked far lower overall than hockey knowledge would justify:
Makar at \#44 overall, for example.

\subsection{Solution: weighted blend with historical average}

For each defense target, the model prediction is blended with a 2-year historical per-game
average that is already present in the engineered feature set:

\begin{equation}
  \hat{y}_i^{\text{blend}} = \alpha \cdot \hat{y}_i^{\text{model}} + (1-\alpha) \cdot \bar{y}_i^{\text{hist}}
  \label{eq:blend}
\end{equation}

If no historical average is available for player $i$ (e.g.\ a new player in the league),
the blend falls back to the raw model prediction.

\subsection{Calibrating alpha from OOS concordance}

The blend weight $\alpha$ is calibrated from out-of-sample concordance. When concordance = 0.5
(chance), the model adds no ranking information and should receive zero weight.
When concordance = 1.0 (perfect ranking), the model should receive full weight.
The natural calibration is:

\begin{equation}
  \alpha = 2 \times (C_{\text{OOS}} - 0.5)
  \label{eq:alpha}
\end{equation}

Applied values for the 2026-27 season:

\begin{center}
\begin{tabular}{lccl}
\toprule
Target & OOS Concordance & $\alpha$ & Interpretation \\
\midrule
\texttt{defense\_points}      & $\approx 0.72$ & 0.45 & Model trusted; hist corrects compression \\
\texttt{defense\_plus\_minus} & $\approx 0.54$ & 0.10 & Model barely beats chance; hist dominates \\
\bottomrule
\end{tabular}
\end{center}

\textbf{Effect:} Makar moved from \#44 overall $\to$ \#4 overall after blending.


% ===========================================================================
\section{Pool Scoring and Ranking}
\label{sec:scoring}
% ===========================================================================

\subsection{Projected games}

Each player's projected games played is derived from a 3-year rolling average of their actual
\texttt{games\_played\_player}, capped at \texttt{cfg.games\_per\_season}.

\textbf{Partial-season filter:} Historical seasons with GP $<$ 25 are excluded from the rolling
window before computing the average. This prevents late-season rookie call-ups (e.g.\ a player
who plays 15 games in their debut year after being recalled from the AHL in March) from biasing
the projection downward. The most-recent season is \emph{always} included regardless of GP —
a genuine recent injury should still reduce the projection.

\begin{equation}
  \mathcal{H}_i = \{\text{GP}_{i,s} : s < t_{\max},\; \text{GP}_{i,s} \geq 25\}
  \quad\text{(filtered historical seasons)}
\end{equation}
\begin{equation}
  \text{GP}_i^{\text{proj}} = \min\!\left(\frac{1}{|\mathcal{W}_i|}\sum_{s \in \mathcal{W}_i}
  \text{GP}_{i,s},\; \texttt{games\_per\_season}\right),
  \quad \mathcal{W}_i = (\mathcal{H}_i \cup \{\text{GP}_{i,t_{\max}}\})\text{ tail-3}
\end{equation}

For players with no qualifying history: forwards and defense default to 70 games; goalies default
to 20 games.

\textbf{GP-projection diagnostic:} On every predict run, a diagnostic compares each player's
projected GP to their most-recent actual GP. Players where
$\text{GP}^{\text{proj}}_i < 0.70 \times \text{GP}^{\text{last}}_i$
are logged as anomalies and written to \texttt{results/diagnostics/gp\_projection\_check.csv}
(all players sorted by ratio, with columns: player\_id, player\_name, position\_group,
projected\_games, last\_season\_gp, ratio). After the partial-season fix, only 13 players
remain flagged — all with genuinely irregular career patterns (multi-season injuries,
suspensions, retirements) rather than call-up artifacts.

\subsection{Dual scoring outputs}

Two pool-point totals are produced for every player:

\begin{itemize}[leftmargin=*]
  \item \texttt{pool\_points}: uses per-player projected GP (injury-risk-adjusted)
  \item \texttt{pool\_points\_full\_season}: uses flat \texttt{games\_per\_season} = 84 for all players
\end{itemize}

Rankings in the output files are sorted by \texttt{pool\_points} by default.

\subsection{Forward scoring}

\begin{equation}
  \text{pool\_points}_i^{\text{fwd}} = \hat{y}_i^{\text{pts/g}} \times \text{GP}_i^{\text{proj}}
\end{equation}

\subsection{Defense scoring}

\begin{align}
  \widehat{\text{pts}}_i &= \hat{y}_i^{\text{pts/g}} \times \text{GP}_i^{\text{proj}} \\
  \widehat{\text{pm}}_i  &= \hat{y}_i^{\text{pm/g}} \times \text{GP}_i^{\text{proj}} \\
  \text{pool\_points}_i^{\text{def}} &= \widehat{\text{pts}}_i + \widehat{\text{pm}}_i
\end{align}

Blended predictions are used for $\hat{y}^{\text{pts/g}}$ and $\hat{y}^{\text{pm/g}}$.

\subsection{Goalie scoring}

Let $W_i = \hat{y}_i^{\text{wins/g}} \times \text{GP}_i^{\text{proj}}$ and
$S_i = \min(\hat{y}_i^{\text{so/g}} \times \text{GP}_i^{\text{proj}},\; W_i)$
(shutouts capped at wins):

\begin{equation}
  \text{pool\_points}_i^{\text{goalie}} = 1 \cdot (W_i - S_i) + 3 \cdot S_i
    + 10 \cdot \mathbf{1}[\text{best GAA among qualified}]
    + 10 \cdot \mathbf{1}[\text{best Sv\% among qualified}]
\end{equation}

A goalie is \emph{qualified} for the bonus if $\text{GP}_i^{\text{proj}} \geq 40$.
The GAA bonus goes to $\arg\min_i \hat{y}_i^{\text{GAA}}$ among qualified goalies;
the save\% bonus to $\arg\max_i \hat{y}_i^{\text{sv\%}}$.


% ===========================================================================
\section{Diagnostics}
\label{sec:diagnostics}
% ===========================================================================

Model diagnostics are generated automatically at the end of every \texttt{--stage predict} run.

\subsection{Outputs}

\begin{description}[leftmargin=4cm, style=nextline]
  \item[\texttt{results/diagnostics/model\_metrics.csv}]
        One row per model. Columns: model, label, algorithm, n\_train, n\_test, test\_r2,
        train\_r2, baseline\_r2, beats\_baseline, overfitting\_ratio, mae, rmse, cv\_mean,
        spearman\_r, concordance, directional\_acc, bias, oos\_r2, oos\_mae, oos\_n,
        oos\_chosen\_model, oos\_spearman\_r, oos\_concordance, oos\_directional\_acc, oos\_bias,
        oos\_top1\_match\_max, oos\_top1\_match\_min.
  \item[\texttt{plots/diagnostics/pva\_<target>.png}]
        Two-panel plot per model: (left) predicted vs.\ actual scatter with $y=x$ reference line
        and $R^2$ annotation; (right) residuals vs.\ predicted (checks heteroskedasticity).
        Generated from \texttt{y\_test}/\texttt{y\_pred\_test} stored in the joblib artifact —
        no retraining required.
  \item[\texttt{plots/diagnostics/rolling\_oos\_all.png}]
        Grid of OOS $R^2$ by holdout year for each model (one subplot per model). Shows whether
        model quality has been stable over time or has degraded in recent seasons.
  \item[\texttt{results/diagnostics/gp\_projection\_check.csv}]
        All players sorted by $\text{GP}^{\text{proj}} / \text{GP}^{\text{last}}$. Players
        below 0.70 are flagged in the log as anomalies requiring manual review before draft day.
        Produced by \texttt{pool\_ranking.\_write\_gp\_diagnostic()}.

  \item[\texttt{results/diagnostics/source\_reconciliation.csv}]
        NHL API vs.\ MoneyPuck agreement on counting stats. Produced by
        \texttt{common/diagnostics/reconcile.py}.
  \item[\texttt{plots/feature\_importance\_<target>.png}]
        Horizontal bar chart of the top 15 features by importance (tree models: built-in
        feature importances; linear models: $|\text{coef}|$).
\end{description}

\subsection{Console summary}

A compact table is printed to stdout at the end of every predict run:

\begin{lstlisting}
=== Model fit summary ===
label             algorithm  test_r2  baseline_r2  beats_baseline  spearman_r  concordance  directional_acc  oos_r2  oos_spearman_r  oos_concordance
Fwd Points/G      gradient_boosting   0.712    0.001     True  0.843  0.912   0.81   0.665   0.812   0.897
Def Points/G      ridge               0.531    0.001     True  0.741  0.869   0.74   0.402   0.711   0.851
...
\end{lstlisting}


% ===========================================================================
\section{Results}
\label{sec:results}
% ===========================================================================

\subsection{Model performance}

Table~\ref{tab:results} summarises out-of-sample performance for all seven production models
as of the first full diagnostic run (2026-07-24). Concordance and Spearman $\rho$ for the OOS
evaluation were not yet computed for all targets at that run; estimated ranges are given where
exact values are unavailable.

\begin{table}[h]
\centering
\caption{OOS model performance — all 7 targets (2026-07-24)}
\label{tab:results}
\begin{tabular}{L{3.2cm} C{2.4cm} C{1.5cm} C{1.5cm} C{1.8cm} l}
\toprule
Target & Algorithm & OOS $R^2$ & OOS Conc. & OOS Spear. $\rho$ & Notes \\
\midrule
Fwd Points/G      & Gradient Boosting & 0.665 & $\approx 0.90$ & $\approx 0.83$ & Healthy \\
Def Points/G      & Ridge / GBM       & $\approx 0.40$ & $\approx 0.72$ & $\approx 0.71$ & Post-blend \\
Def +/$-$/G       & ElasticNet        & 0.031 & $\approx 0.54$ & $\approx 0.10$ & Near-chance; hist dominates \\
Goalie Wins/G     & Ridge             & $\approx 0.10$ & $\approx 0.60$ & $\approx 0.35$ & Team-driven \\
Goalie SO/G       & Ridge             & $\approx 0.05$ & $\approx 0.55$ & $\approx 0.20$ & Rare event \\
Goalie GAA        & Ridge             & $-0.203$       & $\approx 0.52$ & $\approx 0.10$ & Below mean; expected \\
Goalie Sv\%       & Lasso             & $-0.839$       & $\approx 0.51$ & $\approx 0.05$ & Below mean; EB helps \\
\bottomrule
\multicolumn{6}{l}{\small Concordance and Spearman $\rho$ values marked $\approx$ are estimates; exact values in model\_metrics.csv.}
\end{tabular}
\end{table}

\subsection{Top-10 overall ranking (2026-07-24)}

\begin{center}
\begin{tabular}{rll}
\toprule
Overall Rank & Player & Position \\
\midrule
1  & N.\ MacKinnon  & F \\
2  & N.\ Kucherov   & F \\
3  & C.\ McDavid    & F \\
4  & C.\ Makar      & D \\
5  & L.\ Draisaitl  & F \\
6  & D.\ Pastrnak   & F \\
7  & E.\ Bouchard   & D \\
8  & N.\ Suzuki     & F \\
9  & M.\ Necas      & F \\
10 & M.\ Celebrini  & F \\
\bottomrule
\end{tabular}
\end{center}

Cale Makar at \#4 reflects the defense blending fix; prior to blending he ranked \#44.


% ===========================================================================
\section{Known Limitations}
\label{sec:limitations}
% ===========================================================================

\begin{description}[leftmargin=2.5cm, style=nextline]

  \item[Goalie precision stats]
        Season-to-season autocorrelation for GAA and save\% is $r \approx 0.09$ — this is a
        fundamental property of goalie performance data, not a model failure. Published ranges
        (Thomas 2006; Macdonald 2010; Schuckers \& Curro 2013) are $r \approx 0.05$--$0.15$
        for raw Sv\% and $r \approx 0.15$--$0.30$ for GSAx. No model can reliably predict
        who will have the best GAA next season from prior stats alone. The pool bonus only
        requires correctly identifying the argmax/argmin among a small number of qualified
        goalies, which is a weaker requirement than accurate regression.

  \item[Defense plus/minus]
        OOS $R^2 = 0.031$ and concordance $\approx 0.54$. The model barely beats the
        chance baseline. This stat is highly dependent on team context (a plus/minus changes
        by $\pm 1$ whether or not the goalie saves a shot that would not be attributed to any
        on-ice skater). Prediction blending (alpha = 0.10 toward historical averages) partially
        corrects the ranking.

  \item[Traded players — historical team context]
        Players who changed teams during a historical season are assigned the \emph{first}
        team listed in the NHL API's comma-joined \texttt{team\_abbrev} field. Season totals
        (goals, assists, games\_played) are correctly preserved. Only the team-context lag
        features (used by the defense model) may be inaccurate for multi-team seasons.
        Per-team game logs are only available for the most recent season, so a majority-team
        fix across all 18 training seasons would require expanding game-log history.

  \item[MoneyPuck coverage]
        $\sim$5.5\% of players are unmatched in MoneyPuck. These are fringe players with
        very few games played who are genuinely absent from MoneyPuck's records, not join
        failures. They receive NaN for all MoneyPuck features, which become 0 after imputation.

  \item[Observation weighting not implemented]
        All training rows carry equal weight regardless of the player's game count. A backup
        goalie with 3 starts influences the loss function equally to a starter with 60. Planned
        fix: \texttt{sample\_weight = min(GP, 40) / 40} via optional threading in
        \texttt{training.py} (see Future Work).

  \item[Pool scoring rules constant year-to-year]
        No time-varying config is needed; the rules have not changed. Any future rule change
        requires only a config edit.

\end{description}


% ===========================================================================
\section{Future Work}
\label{sec:future}
% ===========================================================================

\begin{description}[leftmargin=3cm, style=nextline]

  \item[GP diagnostic Layer 2: game-log debut date]
        The current partial-season filter (GP $<$ 25) works on the aggregate season count. A
        richer signal is already available: \texttt{nhl\_game\_logs.csv} contains
        \texttt{gameDate} in ISO format; the existing \texttt{get\_player\_game\_log()} endpoint
        in \texttt{nhl\_api.py} can fetch any player/season combination. For any player whose
        debut season has GP $<$ 50, fetching their game log and checking the first game date
        would definitively distinguish a March call-up (first game after February) from an
        October injury (first game in November). Deferred because the threshold filter handles
        the known cases; implement if edge cases emerge.

  \item[Observation weighting]
        Thread \texttt{sample\_weight} through \texttt{\_fit\_one()},
        \texttt{train\_and\_select()}, and \texttt{train\_all\_vs\_exclude\_latest()} in
        \texttt{training.py} as an optional parameter (default None = current behavior).
        In goalie \texttt{build\_xy\_for()}, compute weights as
        $\texttt{min(GP, 40) / 40}$ and return as a 4th element. This down-weights 1--5 start
        goalies who contribute mostly noise to the fit.

  \item[MultiTaskElasticNet for (GAA, save\%)]
        Train a single \texttt{sklearn.linear\_model.MultiTaskElasticNetCV} on the
        $[GAA, \text{save\%}]$ pair to enforce joint feature sparsity: a feature is either
        selected for both targets or neither. The algebraic linkage ($r = -0.84$) means both
        targets should share most of their predictive signal. Only implement if observation
        weighting alone is insufficient.

  \item[Joint D-men model]
        Defenseman goals and assists directly contribute to their plus/minus (a scored goal
        counted in points is also counted as $+1$). Predicting plus\_minus as a residual after
        removing the points contribution would isolate the defensive-zone and special-teams
        components. Alternatively, a multi-output regression joint over (points, plus\_minus)
        could be explored.

  \item[Majority-team for historical traded players]
        Expand per-game log collection from the current season only to all historical seasons
        (one API call per player per season for $\sim$18 years). Assign each player to the
        team they played the most games for. This would make the team-context lag features
        more accurate for the $\sim$5--10\% of historical rows involving mid-season trades.

  \item[Interactive draft-day tool]
        A live tool for draft day: show remaining available players ranked by pool points,
        allow filtering by position, mark picked players. Lowest priority; nothing built yet.

  \item[2027-28 season]
        Copy \texttt{2026-27/config.yaml} $\to$ \texttt{2027-28/config.yaml}, update
        \texttt{season}, \texttt{season\_id}, \texttt{roster\_as\_of\_date},
        \texttt{games\_per\_season} (if the NHL expands again), and any
        \texttt{trade\_overrides}. Run \texttt{python3 run\_season.py --season 2027-28 --stage all}.

\end{description}


% ===========================================================================
% References (inline, no BibTeX required)
% ===========================================================================
\section*{References}

\begin{itemize}[leftmargin=*]
  \item Thomas, A.\ C.\ (2006). ``The impact of puck possession and location on ice hockey
        strategy.'' \emph{Proceedings of the Joint Statistical Meetings (Statistics in Sports)}.
        [Empirical Bayes save\% shrinkage; near-zero YoY autocorrelation for Sv\%.]
  \item Macdonald, B.\ (2010). ``A regression-based adjusted plus-minus statistic for NHL
        players.'' \emph{Journal of Quantitative Analysis in Sports}, 7(3).
  \item Schuckers, M.\ \& Curro, J.\ (2013). ``Total Hockey Rating (THoR): A comprehensive
        statistical rating of National Hockey League forwards and defensemen based upon all
        on-ice events.'' \emph{MIT Sloan Sports Analytics Conference}.
  \item Gramacy, R.\ B., Jensen, S.\ T., \& Taddy, M.\ (2013). ``Estimating player contribution
        in hockey with regularized logistic regression.'' \emph{Journal of Quantitative Analysis
        in Sports}, 9(1).
  \item MoneyPuck.com methodology notes: \url{https://moneypuck.com/about.htm}.
        [GSAx definition; shot danger zone classification.]
\end{itemize}

\end{document}
